\documentclass[11pt,a4paper]{article}

% ---- Packages ----
\usepackage[utf8]{inputenc}
\usepackage[T1]{fontenc}
\usepackage{amsmath,amssymb,amsthm}
\usepackage{mathtools}
\usepackage[margin=1in,headheight=14pt]{geometry}
\usepackage{hyperref}
\usepackage{enumitem}
\usepackage{xcolor}
\usepackage{tcolorbox}
\usepackage{fancyhdr}
\usepackage{booktabs}

% ---- Hyperref ----
\hypersetup{
    colorlinks=true,
    linkcolor=red,
    citecolor=blue,
    urlcolor=blue
}

% ---- Header ----
\pagestyle{fancy}
\fancyhf{}
\fancyhead[L]{\textit{Lesson 7: Hilbert Spaces and Operators}}
\fancyhead[R]{\thepage}
\renewcommand{\headrulewidth}{0.4pt}

% ---- Theorem Environments ----
\theoremstyle{definition}
\newtheorem{definition}{Definition}[section]
\newtheorem{example}{Example}[section]
\newtheorem{exercise}{Exercise}[section]

\theoremstyle{plain}
\newtheorem{theorem}{Theorem}[section]
\newtheorem{lemma}[theorem]{Lemma}
\newtheorem{proposition}[theorem]{Proposition}
\newtheorem{corollary}[theorem]{Corollary}

\theoremstyle{remark}
\newtheorem{remark}{Remark}[section]

% ---- Colored Boxes ----
\tcbuselibrary{theorems,skins,breakable}

\newtcolorbox{keyresult}[1][]{
    colback=green!5!white,
    colframe=green!50!black,
    fonttitle=\bfseries,
    title=#1,
    breakable
}

\newtcolorbox{rlconnection}[1][]{
    colback=blue!5!white,
    colframe=blue!50!black,
    fonttitle=\bfseries,
    title=#1,
    breakable
}

\newtcolorbox{warning}[1][]{
    colback=red!5!white,
    colframe=red!50!black,
    fonttitle=\bfseries,
    title=#1,
    breakable
}

% ---- Operators ----
\newcommand{\R}{\mathbb{R}}
\newcommand{\N}{\mathbb{N}}
\newcommand{\Q}{\mathbb{Q}}
\newcommand{\Z}{\mathbb{Z}}
\newcommand{\C}{\mathbb{C}}
\newcommand{\Prob}{\mathbb{P}}
\newcommand{\E}{\mathbb{E}}
\newcommand{\indicator}{\mathbf{1}}
\DeclareMathOperator{\interior}{int}
\DeclareMathOperator{\cl}{cl}
\DeclareMathOperator{\spn}{span}
\DeclareMathOperator{\conv}{conv}
\DeclareMathOperator{\dom}{dom}
\DeclareMathOperator{\range}{range}
\DeclareMathOperator{\nullspace}{null}
\DeclareMathOperator{\diam}{diam}

% ---- Title ----
\title{\textbf{Lesson 7: Hilbert Spaces and Operators}\\[10pt]
\large Inner Product Spaces, Fundamental Theorems, and Fixed Point Theory\\[5pt]
\normalsize Session 7 of 102 $\mid$ Phase I: Mathematical Foundations $\mid$ 8\,h\\
\normalsize Primary text: Kreyszig, \emph{Introductory Functional Analysis with Applications}, Chapters 3--5}
\author{Curriculum: A Rigorous Learning Path to Multi-Agent Deep RL}
\date{}

\begin{document}

\maketitle

\begin{abstract}
This lesson develops the theory of Hilbert spaces, proves the fundamental theorems
of functional analysis, and establishes the fixed point theorems that underpin
reinforcement learning convergence proofs and game-theoretic equilibrium existence.
We begin with inner product spaces and the geometry they provide---projections,
orthogonal decompositions, and best approximation---then prove the Riesz
Representation Theorem identifying a Hilbert space with its dual. We present
complete proofs of the Hahn--Banach Theorem, the Uniform Boundedness Principle,
the Open Mapping Theorem, and the Closed Graph Theorem. Finally, we prove the
Banach Contraction Mapping Theorem (the engine behind value iteration) and state
the Brouwer, Schauder, and Kakutani fixed point theorems (the foundations for
Nash equilibrium existence). Every theorem has a complete proof or detailed
proof sketch, and RL/ML connections are highlighted throughout.
\end{abstract}

\tableofcontents
\newpage

%% ============================================================================
%% SECTION 1: INTRODUCTION
%% ============================================================================
\section{Introduction}\label{sec:introduction}

The passage from normed spaces to inner product spaces provides the crucial
ingredient of \emph{geometry}: an inner product defines angles, orthogonality,
and projections. These geometric tools are not merely elegant abstractions---they
are the backbone of least-squares methods, kernel machines, and the projection
operators used in temporal-difference learning (LSTD, LSPE).

Complete inner product spaces---\emph{Hilbert spaces}---enjoy the richest
structure among Banach spaces. The Projection Theorem, the Riesz Representation
Theorem, and the spectral theory of operators on Hilbert spaces all depend on
this completeness.

Beyond Hilbert spaces, the four fundamental theorems of functional analysis
(Hahn--Banach, Uniform Boundedness, Open Mapping, Closed Graph) are tools
of immense power for proving existence and continuity results in abstract spaces.

Finally, fixed point theorems translate the abstract structure of function spaces
into concrete algorithmic guarantees:

\begin{itemize}[leftmargin=2em]
    \item The \textbf{Banach Contraction Mapping Theorem} guarantees that
          value iteration converges: the Bellman operator is a
          $\gamma$-contraction on the space of bounded value functions.
    \item The \textbf{Brouwer and Schauder Fixed Point Theorems} guarantee
          existence of fixed points for continuous maps on compact convex sets.
    \item The \textbf{Kakutani Fixed Point Theorem} extends Brouwer to
          set-valued maps, providing the standard proof that every finite
          game has a Nash equilibrium.
\end{itemize}

\begin{rlconnection}[Why This Lesson Matters for RL/ML]
\begin{enumerate}[leftmargin=2em]
    \item \textbf{Projections and least-squares}: Temporal-difference methods
          with linear function approximation (LSTD) project the Bellman
          equation onto a subspace. The Projection Theorem makes this rigorous.
    \item \textbf{Kernel methods}: Reproducing kernel Hilbert spaces (RKHS)
          are Hilbert spaces of functions. The Riesz Representation Theorem
          is the starting point for representer theorems.
    \item \textbf{Value iteration convergence}: The Banach Contraction Mapping
          Theorem, applied to the Bellman operator, gives both existence of
          the optimal value function and the geometric convergence rate
          $\|V_n - V^*\| \leq \gamma^n \|V_0 - V^*\|$.
    \item \textbf{Nash equilibrium existence}: The Kakutani Fixed Point Theorem,
          applied to the best-response correspondence, proves existence of
          Nash equilibria in finite games---the theoretical anchor of MARL.
    \item \textbf{Uniform Boundedness}: Ensures that pointwise-bounded
          families of operators are uniformly bounded, a key tool in proving
          convergence of approximation schemes.
\end{enumerate}
\end{rlconnection}

\medskip
\noindent\textbf{Prerequisites.} Familiarity with metric spaces, normed spaces,
Banach spaces, and bounded linear operators as covered in Lesson~6. Measure
theory from Lessons~1--2 is used in $L^2$ examples.

\medskip
\noindent\textbf{Outline.}
Section~\ref{sec:inner-product} develops inner product spaces and proves the
Cauchy--Schwarz inequality.
Section~\ref{sec:hilbert} introduces Hilbert spaces, orthogonality, and
orthogonal decomposition.
Section~\ref{sec:projection} proves the Projection Theorem and studies
projection operators.
Section~\ref{sec:orthonormal} covers orthonormal systems, Bessel's inequality,
and Parseval's identity.
Section~\ref{sec:riesz} proves the Riesz Representation Theorem.
Section~\ref{sec:fundamental} presents the four fundamental theorems of
functional analysis with complete proofs.
Section~\ref{sec:fixed-point} proves the Banach, Brouwer, Schauder, and
Kakutani fixed point theorems.
Section~\ref{sec:exercises} provides exercises.
Section~\ref{sec:summary} summarizes the key results.


%% ============================================================================
%% SECTION 2: INNER PRODUCT SPACES
%% ============================================================================
\section{Inner Product Spaces}\label{sec:inner-product}

\begin{definition}[Inner Product Space]\label{def:inner-product}
Let $X$ be a vector space over $\mathbb{F} \in \{\R, \C\}$. An
\emph{inner product} on $X$ is a map
$\langle \cdot, \cdot \rangle \colon X \times X \to \mathbb{F}$ satisfying,
for all $x, y, z \in X$ and $\alpha \in \mathbb{F}$:
\begin{enumerate}[label=(\roman*)]
    \item \textbf{Linearity in the first argument:}
          $\langle \alpha x + y, z \rangle = \alpha \langle x, z \rangle
          + \langle y, z \rangle$.
    \item \textbf{Conjugate symmetry:}
          $\langle x, y \rangle = \overline{\langle y, x \rangle}$.
    \item \textbf{Positive definiteness:}
          $\langle x, x \rangle \geq 0$, with equality if and only if $x = 0$.
\end{enumerate}
A vector space equipped with an inner product is called an
\emph{inner product space} (or \emph{pre-Hilbert space}).
\end{definition}

\begin{remark}\label{rem:convention}
We adopt the convention that the inner product is linear in the first argument
and conjugate-linear in the second. Some texts use the opposite convention.
In the real case, conjugate symmetry reduces to symmetry:
$\langle x, y \rangle = \langle y, x \rangle$.
\end{remark}

\begin{example}[Standard Inner Product Spaces]\label{ex:standard-ip}
\begin{enumerate}[label=(\alph*)]
    \item \textbf{Euclidean space $\R^n$:}
          $\langle x, y \rangle = \sum_{i=1}^n x_i y_i$.
    \item \textbf{Sequence space $\ell^2$:}
          $\langle x, y \rangle = \sum_{i=1}^\infty x_i \overline{y_i}$
          for $x = (x_i), y = (y_i) \in \ell^2$.
    \item \textbf{Function space $L^2(\mu)$:}
          $\langle f, g \rangle = \int f \,\overline{g}\, d\mu$
          for $f, g \in L^2(\mu)$.
    \item \textbf{Matrix space $\R^{m \times n}$:}
          $\langle A, B \rangle = \mathrm{tr}(B^\top A)$, the Frobenius
          inner product.
\end{enumerate}
\end{example}

\subsection{The Cauchy--Schwarz Inequality}

The single most important inequality in all of analysis.

\begin{theorem}[Cauchy--Schwarz Inequality]\label{thm:cauchy-schwarz}
Let $X$ be an inner product space. For all $x, y \in X$,
\[
    |\langle x, y \rangle| \leq \|x\| \cdot \|y\|,
\]
where $\|x\| = \langle x, x \rangle^{1/2}$. Equality holds if and only if
$x$ and $y$ are linearly dependent.
\end{theorem}

\begin{proof}
If $y = 0$, both sides are zero and the inequality holds trivially.
Assume $y \neq 0$. For any $\alpha \in \mathbb{F}$,
\[
    0 \leq \langle x - \alpha y, x - \alpha y \rangle
    = \|x\|^2 - \alpha \langle y, x \rangle
      - \overline{\alpha} \langle x, y \rangle
      + |\alpha|^2 \|y\|^2.
\]
Choose $\alpha = \langle x, y \rangle / \|y\|^2$. Then
$\overline{\alpha} = \overline{\langle x, y \rangle} / \|y\|^2
= \langle y, x \rangle / \|y\|^2$, and
\begin{align*}
    0 &\leq \|x\|^2
        - \frac{\langle x, y \rangle}{\|y\|^2} \langle y, x \rangle
        - \frac{\langle y, x \rangle}{\|y\|^2} \langle x, y \rangle
        + \frac{|\langle x, y \rangle|^2}{\|y\|^4} \|y\|^2 \\
      &= \|x\|^2
        - \frac{|\langle x, y \rangle|^2}{\|y\|^2}
        - \frac{|\langle x, y \rangle|^2}{\|y\|^2}
        + \frac{|\langle x, y \rangle|^2}{\|y\|^2} \\
      &= \|x\|^2 - \frac{|\langle x, y \rangle|^2}{\|y\|^2}.
\end{align*}
Rearranging gives $|\langle x, y \rangle|^2 \leq \|x\|^2 \|y\|^2$, and
taking square roots yields the result.

For the equality condition: if $x = \beta y$ for some $\beta \in \mathbb{F}$,
then $|\langle x, y \rangle| = |\beta| \|y\|^2 = \|x\| \|y\|$.
Conversely, equality in the above derivation means
$\langle x - \alpha y, x - \alpha y \rangle = 0$, hence $x = \alpha y$.
\end{proof}

\begin{corollary}[Triangle Inequality]\label{cor:triangle}
$\|x + y\| \leq \|x\| + \|y\|$ for all $x, y \in X$.
\end{corollary}

\begin{proof}
$\|x + y\|^2 = \|x\|^2 + 2\,\mathrm{Re}\,\langle x, y \rangle + \|y\|^2
\leq \|x\|^2 + 2\|x\|\|y\| + \|y\|^2 = (\|x\| + \|y\|)^2$.
\end{proof}

Thus $\|x\| = \langle x, x \rangle^{1/2}$ defines a norm on $X$: the
\emph{induced norm}. Every inner product space is a normed space.

\begin{theorem}[Parallelogram Law]\label{thm:parallelogram}
In an inner product space,
\[
    \|x + y\|^2 + \|x - y\|^2 = 2(\|x\|^2 + \|y\|^2)
    \quad \text{for all } x, y \in X.
\]
\end{theorem}

\begin{proof}
Expand both sides:
\begin{align*}
    \|x + y\|^2 &= \|x\|^2 + 2\,\mathrm{Re}\,\langle x, y \rangle + \|y\|^2, \\
    \|x - y\|^2 &= \|x\|^2 - 2\,\mathrm{Re}\,\langle x, y \rangle + \|y\|^2.
\end{align*}
Adding these gives the result.
\end{proof}

\begin{remark}\label{rem:parallelogram-converse}
The converse holds: a norm satisfying the parallelogram law arises from an inner
product (recovered via the polarization identity). This means, for example,
that $\ell^p$ for $p \neq 2$ is \emph{not} an inner product space.
\end{remark}

\begin{proposition}[Polarization Identity]\label{prop:polarization}
In a real inner product space,
\[
    \langle x, y \rangle
    = \frac{1}{4}\bigl(\|x+y\|^2 - \|x-y\|^2\bigr).
\]
In a complex inner product space,
\[
    \langle x, y \rangle
    = \frac{1}{4}\sum_{k=0}^{3} i^k \|x + i^k y\|^2.
\]
\end{proposition}

\begin{proof}
For the real case, expand $\|x+y\|^2 - \|x-y\|^2 = 4\,\langle x, y \rangle$
using the same computation as in the parallelogram law proof.

For the complex case, let $\omega = i$ and compute:
\begin{align*}
    \sum_{k=0}^{3} i^k \|x + i^k y\|^2
    &= \sum_{k=0}^{3} i^k \bigl(\|x\|^2 + i^k \langle y, x \rangle
       + \overline{i^k} \langle x, y \rangle + \|y\|^2\bigr).
\end{align*}
Since $\sum_{k=0}^{3} i^k = 0$ and $\sum_{k=0}^{3} i^{2k} = 0$,
the $\|x\|^2$ and $\|y\|^2$ terms vanish. We have
$\sum_{k=0}^{3} i^k \cdot \overline{i^k} = \sum_{k=0}^{3} |i^k|^2 = 4$
and $\sum_{k=0}^{3} i^k \cdot i^k = \sum_{k=0}^{3} i^{2k} = 0$.
Therefore the sum equals $4\,\langle x, y \rangle$.
\end{proof}

\begin{proposition}[Continuity of the Inner Product]\label{prop:ip-continuity}
The inner product is jointly continuous: if $x_n \to x$ and $y_n \to y$ in an
inner product space, then $\langle x_n, y_n \rangle \to \langle x, y \rangle$.
\end{proposition}

\begin{proof}
$|\langle x_n, y_n \rangle - \langle x, y \rangle|
\leq |\langle x_n - x, y_n \rangle| + |\langle x, y_n - y \rangle|
\leq \|x_n - x\| \|y_n\| + \|x\| \|y_n - y\|$.
Since convergent sequences are bounded, $\|y_n\| \leq M$ for some $M$,
so the right side tends to zero.
\end{proof}


%% ============================================================================
%% SECTION 3: HILBERT SPACES
%% ============================================================================
\section{Hilbert Spaces}\label{sec:hilbert}

\begin{definition}[Hilbert Space]\label{def:hilbert-space}
A \emph{Hilbert space} is an inner product space that is complete with respect
to the norm induced by its inner product. That is, every Cauchy sequence
converges in the norm $\|x\| = \langle x, x \rangle^{1/2}$.
\end{definition}

\begin{example}[Standard Hilbert Spaces]\label{ex:hilbert-spaces}
\begin{enumerate}[label=(\alph*)]
    \item $\R^n$ with the Euclidean inner product is a Hilbert space (finite-dimensional
          normed spaces are always complete).
    \item $\ell^2 = \{(x_i)_{i=1}^\infty : \sum |x_i|^2 < \infty\}$ is a Hilbert space.
    \item $L^2(\mu) = \{f : \int |f|^2 \, d\mu < \infty\} / {\sim}$ (equivalence
          classes modulo a.e.\ equality) is a Hilbert space by the Riesz--Fischer theorem.
\end{enumerate}
\end{example}

\begin{warning}[Non-Example]
The space $C[0,1]$ with the $L^2$ inner product
$\langle f, g \rangle = \int_0^1 f(t)g(t)\, dt$ is an inner product space
but \emph{not} a Hilbert space: one can construct Cauchy sequences of
continuous functions whose limit in $L^2$ is not continuous.
\end{warning}

\subsection{Orthogonality}

\begin{definition}[Orthogonality]\label{def:orthogonality}
Two elements $x, y$ in an inner product space are \emph{orthogonal}, written
$x \perp y$, if $\langle x, y \rangle = 0$. A set $S \subset X$ is
\emph{orthogonal} if any two distinct elements of $S$ are orthogonal.
\end{definition}

\begin{theorem}[Pythagorean Theorem]\label{thm:pythagorean}
If $x_1, \ldots, x_n$ are pairwise orthogonal, then
\[
    \Bigl\|\sum_{k=1}^n x_k\Bigr\|^2 = \sum_{k=1}^n \|x_k\|^2.
\]
\end{theorem}

\begin{proof}
$\bigl\|\sum_k x_k\bigr\|^2 = \sum_k \sum_j \langle x_k, x_j \rangle
= \sum_k \langle x_k, x_k \rangle = \sum_k \|x_k\|^2$,
since $\langle x_k, x_j \rangle = 0$ for $k \neq j$.
\end{proof}

\begin{definition}[Orthogonal Complement]\label{def:orthogonal-complement}
For a subset $S \subset X$, the \emph{orthogonal complement} of $S$ is
\[
    S^\perp = \{x \in X : \langle x, s \rangle = 0 \text{ for all } s \in S\}.
\]
\end{definition}

\begin{proposition}[Properties of Orthogonal Complements]\label{prop:orth-complement}
Let $S, T$ be subsets of an inner product space $X$. Then:
\begin{enumerate}[label=(\roman*)]
    \item $S^\perp$ is a closed subspace of $X$.
    \item $S \subset T \implies T^\perp \subset S^\perp$.
    \item $S \subset (S^\perp)^\perp$.
    \item $S^\perp = (\overline{\spn(S)})^\perp$.
\end{enumerate}
\end{proposition}

\begin{proof}
\begin{enumerate}[label=(\roman*)]
    \item Linearity of the inner product in the first argument shows $S^\perp$
          is a subspace. Continuity of the inner product
          (Proposition~\ref{prop:ip-continuity}) shows it is closed:
          if $x_n \in S^\perp$ and $x_n \to x$, then for every $s \in S$,
          $\langle x, s \rangle = \lim_n \langle x_n, s \rangle = 0$.
    \item If $x \in T^\perp$, then $\langle x, t \rangle = 0$ for all $t \in T
          \supset S$, so $x \in S^\perp$.
    \item If $s \in S$ and $x \in S^\perp$, then $\langle s, x \rangle = 0$,
          so $s \in (S^\perp)^\perp$.
    \item One inclusion: $S \subset \overline{\spn(S)}$ implies
          $(\overline{\spn(S)})^\perp \subset S^\perp$.
          Other inclusion: if $x \in S^\perp$, then $\langle x, s \rangle = 0$
          for all $s \in S$. By linearity, $\langle x, y \rangle = 0$ for all
          $y \in \spn(S)$. By continuity, $\langle x, y \rangle = 0$ for all
          $y \in \overline{\spn(S)}$. \qedhere
\end{enumerate}
\end{proof}


%% ============================================================================
%% SECTION 4: PROJECTION THEOREM
%% ============================================================================
\section{The Projection Theorem}\label{sec:projection}

\begin{theorem}[Best Approximation in Closed Convex Sets]\label{thm:best-approx}
Let $H$ be a Hilbert space and $C \subset H$ a nonempty, closed, convex set.
For every $x \in H$, there exists a unique $y_0 \in C$ such that
\[
    \|x - y_0\| = \inf_{y \in C} \|x - y\| =: d(x, C).
\]
\end{theorem}

\begin{proof}
Let $d = \inf_{y \in C} \|x - y\|$. Choose a minimizing sequence
$(y_n)$ in $C$ with $\|x - y_n\| \to d$.

\textbf{Claim:} $(y_n)$ is Cauchy. By the parallelogram law applied to
$a = x - y_n$ and $b = x - y_m$:
\[
    \|a + b\|^2 + \|a - b\|^2 = 2(\|a\|^2 + \|b\|^2).
\]
Here $a - b = y_m - y_n$ and
$a + b = 2x - y_n - y_m = 2\bigl(x - \frac{y_n + y_m}{2}\bigr)$.
Since $C$ is convex, $\frac{y_n + y_m}{2} \in C$, so
$\|a + b\| = 2\bigl\|x - \frac{y_n + y_m}{2}\bigr\| \geq 2d$.
Therefore
\[
    \|y_n - y_m\|^2
    = 2(\|x - y_n\|^2 + \|x - y_m\|^2) - 4\Bigl\|x - \frac{y_n + y_m}{2}\Bigr\|^2
    \leq 2(\|x - y_n\|^2 + \|x - y_m\|^2) - 4d^2.
\]
As $n, m \to \infty$, $\|x - y_n\|^2 \to d^2$ and $\|x - y_m\|^2 \to d^2$,
so $\|y_n - y_m\|^2 \to 2d^2 + 2d^2 - 4d^2 = 0$. Thus $(y_n)$ is Cauchy.

Since $H$ is complete and $C$ is closed, $y_n \to y_0 \in C$, and
$\|x - y_0\| = \lim \|x - y_n\| = d$.

\textbf{Uniqueness:} If $y_0, y_0'$ both achieve the infimum, apply the
parallelogram law to $a = x - y_0, b = x - y_0'$:
$\|y_0 - y_0'\|^2 = 2d^2 + 2d^2 - 4\|x - \frac{y_0 + y_0'}{2}\|^2
\leq 4d^2 - 4d^2 = 0$, so $y_0 = y_0'$.
\end{proof}

\begin{theorem}[Projection Theorem]\label{thm:projection}
Let $H$ be a Hilbert space and $M \subset H$ a closed subspace. For every
$x \in H$, there exist unique elements $m \in M$ and $z \in M^\perp$ such that
\[
    x = m + z.
\]
The element $m$ is the unique best approximation of $x$ in $M$, and is
characterized by the condition $x - m \in M^\perp$.
\end{theorem}

\begin{proof}
By Theorem~\ref{thm:best-approx}, there exists a unique $m \in M$ minimizing
$\|x - m\|$. Set $z = x - m$.

\textbf{Claim:} $z \in M^\perp$. For any $v \in M$ and $\alpha \in \mathbb{F}$,
$m + \alpha v \in M$ (since $M$ is a subspace), so
\[
    \|z\|^2 \leq \|x - (m + \alpha v)\|^2 = \|z - \alpha v\|^2
    = \|z\|^2 - 2\,\mathrm{Re}\,(\alpha \langle v, z \rangle)
      + |\alpha|^2 \|v\|^2.
\]
For $v \neq 0$, choose $\alpha = \langle z, v \rangle / \|v\|^2$.
Then $\overline{\alpha} = \langle v, z \rangle / \|v\|^2$ and
$\alpha \langle v, z \rangle = |\langle z, v \rangle|^2 / \|v\|^2$, giving
\[
    0 \leq -\frac{|\langle z, v \rangle|^2}{\|v\|^2},
\]
hence $\langle z, v \rangle = 0$ for all $v \in M$, i.e., $z \in M^\perp$.

\textbf{Uniqueness:} If $x = m_1 + z_1 = m_2 + z_2$ with $m_i \in M$ and
$z_i \in M^\perp$, then $m_1 - m_2 = z_2 - z_1 \in M \cap M^\perp$.
But $\langle w, w \rangle = 0$ for any $w \in M \cap M^\perp$, so $w = 0$.
\end{proof}

\begin{keyresult}[Orthogonal Decomposition]
If $M$ is a closed subspace of a Hilbert space $H$, then
$H = M \oplus M^\perp$ (internal direct sum), meaning every $x \in H$ has
a unique decomposition $x = m + z$ with $m \in M$, $z \in M^\perp$.
Moreover, $(M^\perp)^\perp = M$.
\end{keyresult}

\begin{definition}[Orthogonal Projection]\label{def:projection-operator}
The \emph{orthogonal projection} onto a closed subspace $M$ of a Hilbert
space $H$ is the map $P_M \colon H \to H$ defined by $P_M(x) = m$, where
$x = m + z$ is the unique decomposition from Theorem~\ref{thm:projection}.
\end{definition}

\begin{proposition}[Properties of Orthogonal Projections]\label{prop:projection-props}
Let $P = P_M$ be the orthogonal projection onto a closed subspace $M$. Then:
\begin{enumerate}[label=(\roman*)]
    \item $P$ is linear.
    \item $P^2 = P$ (idempotent).
    \item $\range(P) = M$ and $\nullspace(P) = M^\perp$.
    \item $\langle Px, y \rangle = \langle x, Py \rangle$ for all $x, y \in H$
          (self-adjoint).
    \item $\|P\| = 1$ if $M \neq \{0\}$ (bounded with operator norm 1).
    \item $\|x - Px\| \leq \|x - m\|$ for all $m \in M$, with equality iff
          $m = Px$.
\end{enumerate}
\end{proposition}

\begin{proof}
\begin{enumerate}[label=(\roman*)]
    \item If $x = m_1 + z_1$ and $y = m_2 + z_2$, then
          $\alpha x + \beta y = (\alpha m_1 + \beta m_2) + (\alpha z_1 + \beta z_2)$,
          which is the decomposition since $M$ and $M^\perp$ are subspaces.
          Thus $P(\alpha x + \beta y) = \alpha m_1 + \beta m_2 = \alpha Px + \beta Py$.
    \item $Px \in M$, so $Px = Px + 0$ is the decomposition of $Px$,
          giving $P(Px) = Px$.
    \item Immediate from the decomposition.
    \item $\langle Px, y \rangle = \langle Px, Py + (I - P)y \rangle
          = \langle Px, Py \rangle + \langle Px, (I-P)y \rangle$.
          Since $Px \in M$ and $(I-P)y \in M^\perp$, the second term vanishes.
          Similarly $\langle x, Py \rangle = \langle Px, Py \rangle$.
    \item $\|Px\|^2 = \langle Px, Px \rangle = \langle Px, x \rangle
          \leq \|Px\| \|x\|$ by Cauchy--Schwarz, so $\|Px\| \leq \|x\|$,
          hence $\|P\| \leq 1$. For $0 \neq m \in M$, $Pm = m$, so $\|P\| = 1$.
    \item This is the best approximation property from Theorem~\ref{thm:best-approx}. \qedhere
\end{enumerate}
\end{proof}

\begin{rlconnection}[Projections in Temporal-Difference Learning]
In linear function approximation for RL, the value function is approximated
in a subspace $M = \{X w : w \in \R^d\}$ where $X$ is a feature matrix.
The LSTD (Least-Squares TD) algorithm computes an approximation
$\hat{V} = P_M T V$ where $T$ is the Bellman operator and $P_M$ is the
orthogonal projection onto $M$ (in a weighted $L^2$ norm). Understanding
that $P_M$ is a bounded, idempotent, self-adjoint operator with
$\|P_M\| = 1$ is essential for proving that the projected Bellman equation
$V = P_M T V$ has a unique fixed point and that LSTD converges to it.
\end{rlconnection}


%% ============================================================================
%% SECTION 5: ORTHONORMAL SYSTEMS
%% ============================================================================
\section{Orthonormal Systems}\label{sec:orthonormal}

\begin{definition}[Orthonormal Set]\label{def:orthonormal}
A set $\{e_\alpha\}_{\alpha \in A}$ in an inner product space is
\emph{orthonormal} if
\[
    \langle e_\alpha, e_\beta \rangle = \delta_{\alpha\beta}
    = \begin{cases} 1 & \text{if } \alpha = \beta, \\ 0 & \text{if } \alpha \neq \beta. \end{cases}
\]
\end{definition}

\begin{theorem}[Gram--Schmidt Process]\label{thm:gram-schmidt}
Let $(x_n)_{n=1}^\infty$ be a linearly independent sequence in an inner
product space $X$. Define
\begin{align*}
    e_1 &= \frac{x_1}{\|x_1\|}, \\
    \tilde{e}_n &= x_n - \sum_{k=1}^{n-1} \langle x_n, e_k \rangle e_k,
    \qquad e_n = \frac{\tilde{e}_n}{\|\tilde{e}_n\|}
    \quad (n \geq 2).
\end{align*}
Then $(e_n)$ is an orthonormal sequence with
$\spn\{e_1, \ldots, e_n\} = \spn\{x_1, \ldots, x_n\}$ for every $n$.
\end{theorem}

\begin{proof}
By induction. The base case $n = 1$ is clear. Suppose $e_1, \ldots, e_{n-1}$
are orthonormal with $\spn\{e_1, \ldots, e_{n-1}\} = \spn\{x_1, \ldots, x_{n-1}\}$.
Then $\tilde{e}_n = x_n - \sum_{k=1}^{n-1} \langle x_n, e_k \rangle e_k$
satisfies $\langle \tilde{e}_n, e_j \rangle = \langle x_n, e_j \rangle
- \langle x_n, e_j \rangle = 0$ for $j = 1, \ldots, n-1$.
Since $(x_k)$ is linearly independent, $x_n \notin \spn\{x_1, \ldots, x_{n-1}\}$,
so $\tilde{e}_n \neq 0$. Setting $e_n = \tilde{e}_n / \|\tilde{e}_n\|$ gives
$\|e_n\| = 1$ and $\langle e_n, e_j \rangle = 0$ for $j < n$.
The span equality follows since $x_n = \tilde{e}_n + \sum_{k=1}^{n-1}
\langle x_n, e_k \rangle e_k$ expresses $x_n$ in terms of $e_1, \ldots, e_n$,
and conversely $e_n$ is a linear combination of $x_1, \ldots, x_n$.
\end{proof}

\begin{definition}[Fourier Coefficients]\label{def:fourier}
Given an orthonormal set $\{e_\alpha\}_{\alpha \in A}$ and $x \in X$, the
\emph{Fourier coefficients} of $x$ are $\hat{x}_\alpha = \langle x, e_\alpha \rangle$.
\end{definition}

\begin{theorem}[Bessel's Inequality]\label{thm:bessel}
Let $\{e_n\}_{n=1}^\infty$ be an orthonormal sequence in an inner product
space $X$. For every $x \in X$,
\[
    \sum_{n=1}^\infty |\langle x, e_n \rangle|^2 \leq \|x\|^2.
\]
\end{theorem}

\begin{proof}
For any $N \in \N$, let $s_N = \sum_{n=1}^N \langle x, e_n \rangle e_n$.
Then
\[
    0 \leq \|x - s_N\|^2
    = \|x\|^2 - 2\,\mathrm{Re}\sum_{n=1}^N \langle x, e_n \rangle
      \langle e_n, x \rangle + \sum_{n=1}^N |\langle x, e_n \rangle|^2
    = \|x\|^2 - \sum_{n=1}^N |\langle x, e_n \rangle|^2,
\]
where we used orthonormality to compute $\|s_N\|^2 = \sum_{n=1}^N
|\langle x, e_n \rangle|^2$ (Pythagorean theorem) and expanded the cross term.
Rearranging gives $\sum_{n=1}^N |\langle x, e_n \rangle|^2 \leq \|x\|^2$.
Since this holds for all $N$, the series converges and the inequality persists.
\end{proof}

\begin{definition}[Orthonormal Basis (Hilbert Basis)]\label{def:onb}
An orthonormal set $\{e_\alpha\}_{\alpha \in A}$ in a Hilbert space $H$ is an
\emph{orthonormal basis} (or \emph{complete orthonormal system}) if
$\overline{\spn\{e_\alpha : \alpha \in A\}} = H$.
Equivalently, $\{e_\alpha\}^\perp = \{0\}$.
\end{definition}

\begin{theorem}[Parseval's Identity and Fourier Expansion]\label{thm:parseval}
Let $\{e_n\}_{n=1}^\infty$ be an orthonormal basis of a Hilbert space $H$.
Then for every $x \in H$:
\begin{enumerate}[label=(\roman*)]
    \item \textbf{Fourier expansion:}
          $x = \sum_{n=1}^\infty \langle x, e_n \rangle e_n$
          (convergence in norm).
    \item \textbf{Parseval's identity:}
          $\|x\|^2 = \sum_{n=1}^\infty |\langle x, e_n \rangle|^2$.
    \item \textbf{Generalized Parseval:}
          $\langle x, y \rangle = \sum_{n=1}^\infty \langle x, e_n \rangle
          \overline{\langle y, e_n \rangle}$.
\end{enumerate}
\end{theorem}

\begin{proof}
Let $s_N = \sum_{n=1}^N \langle x, e_n \rangle e_n$.

\textbf{(i)} We show $s_N \to x$. Let $M_N = \spn\{e_1, \ldots, e_N\}$.
By Proposition~\ref{prop:projection-props}, $s_N = P_{M_N} x$ (the projection
onto $M_N$) since $x - s_N \perp e_n$ for $n = 1, \ldots, N$.
For any $\varepsilon > 0$, since $\overline{\spn\{e_n\}} = H$, there exist
$N$ and scalars $c_1, \ldots, c_N$ with $\|x - \sum_{n=1}^N c_n e_n\| < \varepsilon$.
Then $\|x - s_N\| \leq \|x - \sum c_n e_n\| < \varepsilon$ by best
approximation. Thus $\|x - s_N\| \to 0$.

\textbf{(ii)} $\|x\|^2 = \lim_N \|s_N\|^2 + \|x - s_N\|^2$... more directly,
$\|x\|^2 = \lim_N (\|s_N\|^2 + \|x - s_N\|^2)$ by the Pythagorean theorem
(since $s_N \perp (x - s_N)$). Since $\|x - s_N\| \to 0$,
$\|x\|^2 = \lim_N \|s_N\|^2 = \sum_{n=1}^\infty |\langle x, e_n \rangle|^2$.

\textbf{(iii)} Expand:
$\langle x, y \rangle = \langle \sum_n \langle x, e_n \rangle e_n,
\sum_m \langle y, e_m \rangle e_m \rangle$. By continuity of the inner
product and orthonormality,
$= \sum_n \langle x, e_n \rangle \overline{\langle y, e_n \rangle}$.
\end{proof}

\begin{proposition}[Separability and Orthonormal Bases]\label{prop:separability}
A Hilbert space is separable if and only if it admits a countable orthonormal
basis.
\end{proposition}

\begin{proof}
$(\Leftarrow)$ If $\{e_n\}_{n=1}^\infty$ is a countable orthonormal basis,
then finite rational linear combinations of the $e_n$ form a countable
dense subset.

$(\Rightarrow)$ Let $D = \{d_n\}$ be a countable dense subset. Apply
Gram--Schmidt to any maximal linearly independent subsequence. The resulting
orthonormal sequence has the same closed span as $D$, which is $H$.
\end{proof}


%% ============================================================================
%% SECTION 6: RIESZ REPRESENTATION THEOREM
%% ============================================================================
\section{The Riesz Representation Theorem}\label{sec:riesz}

\begin{theorem}[Riesz Representation Theorem]\label{thm:riesz}
Let $H$ be a Hilbert space and $f \colon H \to \mathbb{F}$ a bounded linear
functional ($f \in H^*$). Then there exists a unique $y \in H$ such that
\[
    f(x) = \langle x, y \rangle \quad \text{for all } x \in H.
\]
Moreover, $\|f\|_{H^*} = \|y\|_H$.
\end{theorem}

\begin{proof}
\textbf{Existence.} If $f = 0$, take $y = 0$.
Suppose $f \neq 0$. Let $M = \ker f = \{x \in H : f(x) = 0\}$, which is
a closed subspace (since $f$ is continuous). Since $f \neq 0$, there exists
$z \in M^\perp$ with $z \neq 0$.

By the Projection Theorem, $H = M \oplus M^\perp$. We claim
$\dim(M^\perp) = 1$: if $z_1, z_2 \in M^\perp \setminus \{0\}$, then
$f(z_2) z_1 - f(z_1) z_2 \in M$ (check: $f(f(z_2)z_1 - f(z_1)z_2) = 0$)
and also in $M^\perp$, so it equals 0. Thus $z_1$ and $z_2$ are linearly
dependent, confirming $\dim(M^\perp) = 1$.

Hence $M^\perp = \spn\{z\}$ for some $z \neq 0$ with $f(z) \neq 0$.
Every $x \in H$ can be written $x = m + \alpha z$ with $m \in M$, so
$f(x) = \alpha f(z)$.

Also, $\langle x, z \rangle = \langle m + \alpha z, z \rangle
= \alpha \|z\|^2$ (since $m \perp z$), giving $\alpha = \langle x, z \rangle / \|z\|^2$.

Therefore
$f(x) = \frac{f(z)}{\|z\|^2} \langle x, z \rangle
= \langle x, \frac{\overline{f(z)}}{\|z\|^2} z \rangle$.
Set $y = \frac{\overline{f(z)}}{\|z\|^2} z$.

\textbf{Uniqueness.} If $\langle x, y_1 \rangle = \langle x, y_2 \rangle$
for all $x$, then $\langle x, y_1 - y_2 \rangle = 0$ for all $x$.
Taking $x = y_1 - y_2$ gives $\|y_1 - y_2\|^2 = 0$.

\textbf{Norm equality.}
$|f(x)| = |\langle x, y \rangle| \leq \|x\| \|y\|$ by Cauchy--Schwarz,
so $\|f\| \leq \|y\|$. Taking $x = y$:
$f(y) = \langle y, y \rangle = \|y\|^2$, so $\|f\| \geq |f(y/\|y\|)| = \|y\|$.
\end{proof}

\begin{keyresult}[Riesz Representation -- Dual of a Hilbert Space]
The map $\Phi \colon H \to H^*$ defined by $\Phi(y)(x) = \langle x, y \rangle$
is a conjugate-linear isometric isomorphism. Thus $H \cong H^*$: a Hilbert
space is isometrically isomorphic to its dual. This is a property
\emph{unique} to Hilbert spaces among Banach spaces.
\end{keyresult}

\begin{rlconnection}[Riesz Representation and Kernel Methods]
In a reproducing kernel Hilbert space (RKHS) $\mathcal{H}_K$, the evaluation
functional $f \mapsto f(x)$ is bounded linear. By the Riesz Representation
Theorem, there exists $K_x \in \mathcal{H}_K$ with $f(x) = \langle f, K_x \rangle$.
The function $K(x, y) = \langle K_x, K_y \rangle$ is the reproducing kernel.
This is the foundation of kernel methods in ML: support vector machines,
Gaussian processes, and kernel-based policy evaluation in RL all rest on
this identification.
\end{rlconnection}


%% ============================================================================
%% SECTION 7: FUNDAMENTAL THEOREMS OF FUNCTIONAL ANALYSIS
%% ============================================================================
\section{Fundamental Theorems of Functional Analysis}\label{sec:fundamental}

We now present the four pillars of linear functional analysis. These theorems
hold in general Banach spaces (not just Hilbert spaces).

\subsection{Hahn--Banach Theorem}

\begin{definition}[Sublinear Functional]\label{def:sublinear}
A function $p \colon X \to \R$ on a real vector space $X$ is \emph{sublinear} if:
\begin{enumerate}[label=(\roman*)]
    \item $p(\alpha x) = \alpha\, p(x)$ for all $\alpha > 0$ and $x \in X$
          (positive homogeneity).
    \item $p(x + y) \leq p(x) + p(y)$ for all $x, y \in X$ (subadditivity).
\end{enumerate}
\end{definition}

\begin{theorem}[Hahn--Banach Theorem, Real Case]\label{thm:hahn-banach}
Let $X$ be a real vector space, $p \colon X \to \R$ a sublinear functional,
and $Y \subset X$ a subspace. If $f \colon Y \to \R$ is a linear functional
with $f(y) \leq p(y)$ for all $y \in Y$, then there exists a linear functional
$F \colon X \to \R$ extending $f$ (i.e., $F|_Y = f$) such that
$F(x) \leq p(x)$ for all $x \in X$.
\end{theorem}

\begin{proof}
\textbf{Step 1: Extension by one dimension.}
Let $x_0 \in X \setminus Y$ and $Y_1 = Y + \R x_0 = \{y + t x_0 : y \in Y, t \in \R\}$.
We seek $c = F(x_0) \in \R$ such that $F(y + t x_0) = f(y) + tc$ satisfies
$F \leq p$ on $Y_1$.

For $t > 0$: $f(y) + tc \leq p(y + tx_0)$, i.e.,
$c \leq \frac{1}{t}[p(y + tx_0) - f(y)] = p(y/t + x_0) - f(y/t)$.
Since $y/t$ ranges over $Y$, we need
$c \leq \inf_{u \in Y}[p(u + x_0) - f(u)]$.

For $t < 0$: setting $s = -t > 0$, we need $f(y) - sc \leq p(y - sx_0)$,
i.e., $c \geq f(y/s) - p(y/s - x_0)$, so
$c \geq \sup_{v \in Y}[f(v) - p(v - x_0)]$.

We must verify:
$f(v) - p(v - x_0) \leq p(u + x_0) - f(u)$ for all $u, v \in Y$.
This reduces to $f(u) + f(v) \leq p(u + x_0) + p(v - x_0)$.
Indeed, $f(u) + f(v) = f(u + v) \leq p(u + v) = p((u + x_0) + (v - x_0))
\leq p(u + x_0) + p(v - x_0)$.

So we can choose $c$ in the interval
$[\sup_v(f(v) - p(v - x_0)),\, \inf_u(p(u + x_0) - f(u))]$.

\textbf{Step 2: Zorn's Lemma.}
Let $\mathcal{F}$ be the set of all pairs $(Z, g)$ where $Y \subset Z \subset X$
is a subspace and $g \colon Z \to \R$ is linear with $g|_Y = f$ and
$g(z) \leq p(z)$ for all $z \in Z$. Order $\mathcal{F}$ by extension:
$(Z_1, g_1) \leq (Z_2, g_2)$ if $Z_1 \subset Z_2$ and $g_2|_{Z_1} = g_1$.

Every chain has an upper bound (take the union of subspaces and the common
extension). By Zorn's Lemma, $\mathcal{F}$ has a maximal element $(Z^*, F)$.
If $Z^* \neq X$, Step~1 would produce a proper extension, contradicting
maximality. Hence $Z^* = X$.
\end{proof}

\begin{corollary}[Extension of Bounded Linear Functionals]\label{cor:hahn-banach-normed}
Let $Y$ be a subspace of a normed space $X$ and $f \in Y^*$. Then there exists
$F \in X^*$ with $F|_Y = f$ and $\|F\|_{X^*} = \|f\|_{Y^*}$.
\end{corollary}

\begin{proof}
Apply Theorem~\ref{thm:hahn-banach} with $p(x) = \|f\| \cdot \|x\|$.
Then $f(y) \leq |f(y)| \leq \|f\| \|y\| = p(y)$, and the extension $F$
satisfies $F(x) \leq \|f\| \|x\|$ and $F(-x) \leq \|f\| \|x\|$,
giving $|F(x)| \leq \|f\| \|x\|$, so $\|F\| \leq \|f\|$.
Since $F$ extends $f$, $\|F\| \geq \|f\|$.
\end{proof}

\begin{corollary}\label{cor:hahn-banach-separation}
For any $x_0 \neq 0$ in a normed space $X$, there exists $f \in X^*$ with
$\|f\| = 1$ and $f(x_0) = \|x_0\|$.
\end{corollary}

\begin{proof}
Define $g \colon \R x_0 \to \R$ by $g(\alpha x_0) = \alpha \|x_0\|$.
Then $\|g\| = 1$ and $g(x_0) = \|x_0\|$.
Extend by Corollary~\ref{cor:hahn-banach-normed}.
\end{proof}

\subsection{Baire Category Theorem and Uniform Boundedness}

We need a preliminary result.

\begin{theorem}[Baire Category Theorem]\label{thm:baire}
Let $(X, d)$ be a complete metric space. If $X = \bigcup_{n=1}^\infty F_n$
where each $F_n$ is closed, then at least one $F_n$ has nonempty interior.
Equivalently, a countable intersection of open dense sets is dense.
\end{theorem}

\begin{proof}
Suppose for contradiction that every $F_n$ has empty interior. Then each
$G_n = X \setminus F_n$ is open and dense. We construct a nested sequence of
closed balls: since $G_1$ is dense, pick $x_1 \in G_1$ and $r_1 < 1$ with
$\overline{B(x_1, r_1)} \subset G_1$. Since $G_2$ is dense and
$B(x_1, r_1)$ is open, pick $x_2 \in G_2 \cap B(x_1, r_1)$ and $r_2 < 1/2$
with $\overline{B(x_2, r_2)} \subset G_2 \cap B(x_1, r_1)$.

Inductively, we obtain $\overline{B(x_{n+1}, r_{n+1})} \subset G_{n+1}
\cap B(x_n, r_n)$ with $r_n < 1/n$. The centers $(x_n)$ form a Cauchy
sequence (since $d(x_m, x_n) < 2/\min(m,n)$ for $m, n$ large), converging
to some $x^* \in X$ by completeness. Then $x^* \in \overline{B(x_n, r_n)}
\subset G_n$ for all $n$, so $x^* \in \bigcap G_n = X \setminus \bigcup F_n
= \emptyset$, a contradiction.
\end{proof}

\begin{theorem}[Uniform Boundedness Principle / Banach--Steinhaus]\label{thm:ubp}
Let $X$ be a Banach space, $Y$ a normed space, and
$\{T_\alpha\}_{\alpha \in A} \subset \mathcal{B}(X, Y)$ a family of bounded
linear operators. If $\sup_\alpha \|T_\alpha x\| < \infty$ for every $x \in X$
(pointwise bounded), then $\sup_\alpha \|T_\alpha\| < \infty$ (uniformly bounded).
\end{theorem}

\begin{proof}
For each $n \in \N$, define the closed set
\[
    F_n = \{x \in X : \|T_\alpha x\| \leq n \text{ for all } \alpha \in A\}
    = \bigcap_{\alpha \in A} \{x : \|T_\alpha x\| \leq n\}.
\]
Each $F_n$ is closed (as an intersection of closed sets, since each
$x \mapsto \|T_\alpha x\|$ is continuous). Since the family is pointwise
bounded, for each $x$ there exists $n$ with $x \in F_n$, so
$X = \bigcup_{n=1}^\infty F_n$.

By the Baire Category Theorem~\ref{thm:baire}, some $F_N$ has nonempty
interior: there exist $x_0 \in X$ and $r > 0$ with $B(x_0, r) \subset F_N$.

For any $x \in X$ with $\|x\| \leq 1$, we have $x_0 + rx \in B(x_0, r) \subset F_N$
and $x_0 \in F_N$, so for all $\alpha$:
\[
    \|T_\alpha(rx)\| = \|T_\alpha(x_0 + rx) - T_\alpha(x_0)\|
    \leq \|T_\alpha(x_0 + rx)\| + \|T_\alpha(x_0)\| \leq N + N = 2N.
\]
Thus $\|T_\alpha x\| \leq 2N/r$ for all $\|x\| \leq 1$, giving
$\|T_\alpha\| \leq 2N/r$ for all $\alpha$.
\end{proof}

\begin{rlconnection}[Uniform Boundedness in Approximation Theory]
The Uniform Boundedness Principle is used to prove negative results in
approximation theory: certain approximation schemes \emph{must} diverge
for some continuous functions. In RL, it guarantees that if a sequence
of operators (e.g., truncated Bellman operators or finite-sample
approximations) converges pointwise, the operator norms remain controlled.
This is crucial for interchanging limits in convergence proofs.
\end{rlconnection}

\subsection{Open Mapping Theorem}

\begin{theorem}[Open Mapping Theorem]\label{thm:open-mapping}
Let $X$ and $Y$ be Banach spaces and $T \colon X \to Y$ a bounded linear
operator that is surjective. Then $T$ is an open map: $T(U)$ is open in $Y$
whenever $U$ is open in $X$.
\end{theorem}

\begin{proof}
\textbf{Step 1:} We show $T(B_X(0, 1))$ contains a ball centered at 0 in $Y$.
Write $B_X = B_X(0,1)$ and $B_Y = B_Y(0,1)$.

Since $T$ is surjective, $Y = \bigcup_{n=1}^\infty T(nB_X) = \bigcup_{n=1}^\infty n \cdot T(B_X)$.
By Baire's theorem, some $\overline{n \cdot T(B_X)}$ has nonempty interior.
By scaling, $\overline{T(B_X)}$ has nonempty interior: there exist
$y_0 \in Y$ and $s > 0$ with $B_Y(y_0, s) \subset \overline{T(B_X)}$.

Since $\overline{T(B_X)}$ is symmetric (if $y \in \overline{T(B_X)}$, then
$-y \in \overline{T(B_X)}$) and convex, we get
$B_Y(0, s) \subset \overline{T(B_X)}$: for any $\|y\| < s$,
$y + y_0 \in \overline{T(B_X)}$ and $-y + y_0 \in \overline{T(B_X)}$,
and $y_0 \in \overline{T(B_X)}$, so
$y = (y + y_0)/2 + (-y_0 + y)/2 + y_0 - y_0$... Let us be more careful.
We have $y + y_0$ and $-y + y_0$ both in $\overline{T(B_X)}$, and
$\overline{T(B_X)} - \overline{T(B_X)} \subset \overline{T(2B_X)}$
(since $T(B_X) + T(B_X) \subset T(2B_X)$). In particular,
$(y + y_0) - (-y + y_0) = 2y \in \overline{T(2B_X)}$, so
$y \in \overline{T(B_X)}$. Hence $B_Y(0, s) \subset \overline{T(B_X)}$.

\textbf{Step 2:} We show $\overline{T(B_X)} \subset T(2B_X)$, i.e.,
the closure can be ``opened.'' Take $y \in \overline{T(B_X)}$. From Step 1,
$B_Y(0, s) \subset \overline{T(B_X)}$, so for each $r > 0$,
$B_Y(0, rs) \subset \overline{T(rB_X)}$.

Choose $x_1 \in B_X$ with $\|y - Tx_1\| < s/2$. Then
$y - Tx_1 \in B_Y(0, s/2) \subset \overline{T(\frac{1}{2}B_X)}$, so choose
$x_2 \in \frac{1}{2}B_X$ with $\|y - Tx_1 - Tx_2\| < s/4$.
Inductively, choose $x_n \in 2^{-(n-1)} B_X$ with
$\|y - \sum_{k=1}^n Tx_k\| < s \cdot 2^{-n}$.

Then $\sum_n \|x_n\| < 1 + 1/2 + 1/4 + \cdots = 2$. Since $X$ is complete,
$x = \sum_{n=1}^\infty x_n$ converges with $\|x\| < 2$, and
$Tx = \lim_{n \to \infty} T(\sum_{k=1}^n x_k) = y$. So $y \in T(2B_X)$.

\textbf{Step 3:} Combining Steps 1 and 2:
$B_Y(0, s) \subset \overline{T(B_X)} \subset T(2B_X)$, so
$B_Y(0, s/2) \subset T(B_X)$.

\textbf{Step 4:} $T$ is open. Let $U \subset X$ be open and $y_0 = Tx_0 \in T(U)$.
Choose $r > 0$ with $B_X(x_0, r) \subset U$. Then
$T(U) \supset T(B_X(x_0, r)) = Tx_0 + T(rB_X) \supset y_0 + B_Y(0, rs/2)$,
so $T(U)$ is open.
\end{proof}

\begin{corollary}[Bounded Inverse Theorem]\label{cor:bounded-inverse}
If $T \colon X \to Y$ is a bounded linear bijection between Banach spaces,
then $T^{-1}$ is bounded.
\end{corollary}

\begin{proof}
$T$ is an open map, so $T^{-1}$ is continuous (preimages of open sets under
$T^{-1}$ are images of open sets under $T$).
\end{proof}

\subsection{Closed Graph Theorem}

\begin{definition}[Closed Graph]\label{def:closed-graph}
A linear operator $T \colon X \to Y$ between normed spaces has a
\emph{closed graph} if its graph
$G(T) = \{(x, Tx) : x \in X\} \subset X \times Y$ is closed in the
product topology (equivalently, in the norm $\|(x,y)\| = \|x\| + \|y\|$).
Equivalently: if $x_n \to x$ in $X$ and $Tx_n \to y$ in $Y$, then $y = Tx$.
\end{definition}

\begin{theorem}[Closed Graph Theorem]\label{thm:closed-graph}
Let $X$ and $Y$ be Banach spaces. A linear operator $T \colon X \to Y$ is
bounded if and only if it has a closed graph.
\end{theorem}

\begin{proof}
$(\Rightarrow)$ If $T$ is bounded and $x_n \to x$, $Tx_n \to y$, then
$Tx_n \to Tx$ by continuity, so $y = Tx$ by uniqueness of limits.

$(\Leftarrow)$ Assume $G(T)$ is closed. Consider $X \times Y$ with the norm
$\|(x,y)\| = \|x\| + \|y\|$, which is a Banach space. $G(T)$ is a closed
subspace, hence itself a Banach space.

Define $\pi_1 \colon G(T) \to X$ by $\pi_1(x, Tx) = x$ and
$\pi_2 \colon G(T) \to Y$ by $\pi_2(x, Tx) = Tx$.
Then $\pi_1$ is bounded, linear, and bijective. By the Bounded Inverse
Theorem (Corollary~\ref{cor:bounded-inverse}), $\pi_1^{-1} \colon X \to G(T)$
is bounded. Since $T = \pi_2 \circ \pi_1^{-1}$ and both $\pi_2$ and
$\pi_1^{-1}$ are bounded, $T$ is bounded.
\end{proof}

\begin{remark}\label{rem:closed-graph-usage}
The Closed Graph Theorem is often the most convenient way to prove that
an operator is bounded: instead of estimating $\|Tx\|$ directly, one
only needs to verify the sequential closedness condition, which is often
easier in applications.
\end{remark}


%% ============================================================================
%% SECTION 8: FIXED POINT THEOREMS
%% ============================================================================
\section{Fixed Point Theorems}\label{sec:fixed-point}

Fixed point theorems guarantee the existence (and sometimes uniqueness) of
solutions to equations of the form $T(x) = x$. They are among the most
powerful tools in analysis, with direct applications to existence and
convergence theorems in RL.

\subsection{Banach Contraction Mapping Theorem}

\begin{definition}[Contraction]\label{def:contraction}
Let $(X, d)$ be a metric space. A map $T \colon X \to X$ is a
\emph{contraction} (or \emph{contraction mapping}) if there exists
$\gamma \in [0, 1)$ such that
\[
    d(Tx, Ty) \leq \gamma \, d(x, y) \quad \text{for all } x, y \in X.
\]
The constant $\gamma$ is called the \emph{contraction factor}.
\end{definition}

\begin{theorem}[Banach Contraction Mapping Theorem]\label{thm:banach-contraction}
Let $(X, d)$ be a nonempty complete metric space and $T \colon X \to X$ a
contraction with factor $\gamma \in [0,1)$. Then:
\begin{enumerate}[label=(\roman*)]
    \item $T$ has a unique fixed point $x^* \in X$ (i.e., $Tx^* = x^*$).
    \item For any $x_0 \in X$, the iterates $x_{n+1} = Tx_n$ satisfy
          $x_n \to x^*$.
    \item The rate of convergence is geometric:
          $d(x_n, x^*) \leq \frac{\gamma^n}{1 - \gamma} d(x_0, x_1)$.
\end{enumerate}
\end{theorem}

\begin{proof}
\textbf{(i) and (ii): Existence and convergence.}
Fix $x_0 \in X$ and define $x_{n+1} = Tx_n$. First,
\[
    d(x_{n+1}, x_n) = d(Tx_n, Tx_{n-1}) \leq \gamma \, d(x_n, x_{n-1})
    \leq \cdots \leq \gamma^n d(x_1, x_0).
\]
For $m > n$, by the triangle inequality:
\begin{align*}
    d(x_m, x_n) &\leq \sum_{k=n}^{m-1} d(x_{k+1}, x_k)
    \leq \sum_{k=n}^{m-1} \gamma^k d(x_1, x_0)
    \leq \frac{\gamma^n}{1 - \gamma} d(x_1, x_0).
\end{align*}
Since $\gamma < 1$, $\gamma^n \to 0$, so $(x_n)$ is Cauchy. By completeness,
$x_n \to x^*$ for some $x^* \in X$.

Since $T$ is continuous (in fact, Lipschitz), $Tx^* = T(\lim x_n) = \lim Tx_n
= \lim x_{n+1} = x^*$.

\textbf{Uniqueness.} If $Tx^* = x^*$ and $Ty^* = y^*$, then
$d(x^*, y^*) = d(Tx^*, Ty^*) \leq \gamma \, d(x^*, y^*)$, so
$(1 - \gamma) d(x^*, y^*) \leq 0$, giving $d(x^*, y^*) = 0$.

\textbf{(iii): Rate of convergence.}
Letting $m \to \infty$ in the estimate above:
$d(x^*, x_n) = \lim_{m \to \infty} d(x_m, x_n)
\leq \frac{\gamma^n}{1 - \gamma} d(x_1, x_0)$.
\end{proof}

\begin{keyresult}[Banach Contraction Mapping Theorem]
In a complete metric space, a contraction has a unique fixed point, iterative
computation converges from any starting point, and the convergence rate is
geometric: $d(x_n, x^*) \leq \frac{\gamma^n}{1 - \gamma} d(x_0, Tx_0)$.
\end{keyresult}

\begin{rlconnection}[Value Iteration Convergence]
Let $\mathcal{S}$ be a finite state space and consider the Banach space
$(V(\mathcal{S}), \|\cdot\|_\infty)$ of bounded value functions. The
\textbf{Bellman optimality operator} $T \colon V(\mathcal{S}) \to V(\mathcal{S})$
is defined by
\[
    (TV)(s) = \max_{a \in \mathcal{A}} \Bigl[R(s,a) + \gamma \sum_{s'} P(s'|s,a) V(s')\Bigr].
\]
One verifies that $T$ is a $\gamma$-contraction in $\|\cdot\|_\infty$:
\[
    \|TV - TW\|_\infty \leq \gamma \|V - W\|_\infty.
\]
By the Banach Contraction Mapping Theorem:
\begin{enumerate}[leftmargin=2em]
    \item There exists a unique optimal value function $V^*$ with $TV^* = V^*$
          (the Bellman equation).
    \item \textbf{Value iteration} $V_{n+1} = TV_n$ converges to $V^*$ from
          any initialization.
    \item The convergence rate is
          $\|V_n - V^*\|_\infty \leq \frac{\gamma^n}{1 - \gamma} \|V_1 - V_0\|_\infty$.
\end{enumerate}
This is the \emph{fundamental convergence guarantee} for value iteration.
The same argument applies to the Bellman evaluation operator for policy
evaluation, and to $Q$-value iteration.
\end{rlconnection}

\subsection{Brouwer Fixed Point Theorem}

\begin{theorem}[Brouwer Fixed Point Theorem]\label{thm:brouwer}
Let $C \subset \R^n$ be a nonempty compact convex set and $f \colon C \to C$
a continuous map. Then $f$ has a fixed point: there exists $x^* \in C$ with
$f(x^*) = x^*$.
\end{theorem}

\begin{proof}
We prove the theorem for $C = \overline{B}^n = \{x \in \R^n : \|x\| \leq 1\}$
(the closed unit ball). The general case follows since any nonempty compact
convex set in $\R^n$ is homeomorphic to $\overline{B}^n$ (or a lower-dimensional
ball, where the argument is analogous).

\textbf{Proof by contradiction using the retraction argument.}
Suppose $f$ has no fixed point, i.e., $f(x) \neq x$ for all $x \in \overline{B}^n$.
Define $r \colon \overline{B}^n \to S^{n-1}$ (the boundary sphere) as follows:
for each $x$, draw the ray from $f(x)$ through $x$; let $r(x)$ be the point
where this ray hits $S^{n-1}$.

Explicitly, $r(x) = x + t(x)(x - f(x))$ where $t(x) \geq 0$ is chosen so
that $\|r(x)\| = 1$. This leads to a quadratic equation in $t$ which has a
unique nonnegative solution depending continuously on $x$ (since $f(x) \neq x$).

The map $r$ is continuous, and for $x \in S^{n-1}$, we have
$\|x\| = 1$ and the ray from $f(x)$ through $x$ hits $S^{n-1}$ at $x$ itself
(taking $t = 0$), so $r(x) = x$.

Thus $r \colon \overline{B}^n \to S^{n-1}$ is a \emph{retraction}: a
continuous map that is the identity on $S^{n-1}$.

\textbf{Claim:} No such retraction exists. This is the
``no-retraction theorem,'' which can be proved via:
\begin{enumerate}[leftmargin=2em]
    \item \emph{Algebraic topology:} $r$ would induce a surjection
          $r_* \colon \pi_{n-1}(\overline{B}^n) \to \pi_{n-1}(S^{n-1})$,
          but $\pi_{n-1}(\overline{B}^n) = 0$ (the ball is contractible)
          while $\pi_{n-1}(S^{n-1}) \cong \Z$ (for $n \geq 2$), contradiction.
    \item \emph{Smooth approximation + degree theory:} Approximate $r$ by a
          smooth map; the degree of the restriction $r|_{S^{n-1}} = \mathrm{id}$
          is 1, but the degree of any map extending to the ball must be 0.
    \item \emph{Combinatorial (Sperner's lemma):} Triangulate $\overline{B}^n$
          and use Sperner's lemma to show that any continuous map satisfying
          the boundary condition must have a ``completely labeled'' simplex,
          from which a fixed point is obtained in the limit of refinements.
\end{enumerate}
Each approach yields the contradiction. Hence the assumption that $f$ has no
fixed point is false.
\end{proof}

\begin{remark}\label{rem:brouwer-nonconstructive}
Brouwer's theorem is a pure existence result: it gives no algorithm for
finding the fixed point (in fact, computing Brouwer fixed points is
PPAD-complete). For algorithmic purposes, the Banach contraction theorem
is far more useful when it applies.
\end{remark}

\subsection{Schauder Fixed Point Theorem}

The Schauder theorem generalizes Brouwer's theorem to infinite-dimensional
spaces. The key idea is to approximate a compact operator by finite-dimensional
maps and apply Brouwer.

\begin{definition}[Compact Operator]\label{def:compact-op}
A map $T \colon C \to X$ (where $C \subset X$) is \emph{compact} if $T$ is
continuous and maps bounded sets to relatively compact sets (i.e., sets whose
closure is compact).
\end{definition}

\begin{theorem}[Schauder Fixed Point Theorem]\label{thm:schauder}
Let $X$ be a Banach space, $C \subset X$ a nonempty closed bounded convex set,
and $T \colon C \to C$ a compact map. Then $T$ has a fixed point.
\end{theorem}

\begin{proof}
Since $T(C)$ is relatively compact (as $C$ is bounded and $T$ is compact),
$K = \overline{T(C)}$ is compact. For each $n \in \N$, cover $K$ by finitely
many balls of radius $1/n$: there exist $y_1^n, \ldots, y_{k_n}^n \in T(C)$
with $K \subset \bigcup_{i=1}^{k_n} B(y_i^n, 1/n)$.

Define $C_n = \conv\{y_1^n, \ldots, y_{k_n}^n\} \subset C$ (since $C$ is
convex and $y_i^n \in T(C) \subset C$). This is a compact convex subset of
a finite-dimensional subspace.

Define $P_n \colon K \to C_n$ by
\[
    P_n(x) = \frac{\sum_{i=1}^{k_n} \lambda_i^n(x)\, y_i^n}{\sum_{i=1}^{k_n} \lambda_i^n(x)},
    \quad \text{where } \lambda_i^n(x) = \max\bigl(0, 1/n - \|x - y_i^n\|\bigr).
\]
This is the \emph{Schauder projection}. Note that $\sum_i \lambda_i^n(x) > 0$
for $x \in K$ (since $K \subset \bigcup_i B(y_i^n, 1/n)$), so $P_n$ is
well-defined. Moreover, $P_n$ is continuous, maps into $C_n$, and satisfies
\[
    \|P_n(x) - x\| \leq 1/n \quad \text{for all } x \in K.
\]
To see this: $\lambda_i^n(x) > 0$ only when $\|x - y_i^n\| < 1/n$, so
$P_n(x)$ is a convex combination of points within distance $1/n$ of $x$,
hence $\|P_n(x) - x\| < 1/n$.

Now consider $T_n = P_n \circ T|_{C_n} \colon C_n \to C_n$. This is a
continuous self-map of a compact convex subset of a finite-dimensional space.
By the Brouwer Fixed Point Theorem, there exists $x_n \in C_n$ with
$T_n(x_n) = x_n$, i.e.,
\[
    x_n = P_n(Tx_n).
\]
Then $\|x_n - Tx_n\| = \|P_n(Tx_n) - Tx_n\| < 1/n$.

Since $(x_n)$ is a sequence in $C$ and $T$ is compact, $(Tx_n)$ has a
convergent subsequence. Along this subsequence, $x_n - Tx_n \to 0$ and
$Tx_n \to y$, so $x_n \to y$ as well. Since $C$ is closed, $y \in C$.
By continuity of $T$: $Ty = \lim T(x_n) = y$.
\end{proof}

\begin{remark}\label{rem:schauder-variants}
There are several variants of Schauder's theorem. The version above
(sometimes called the Schauder--Tychonoff theorem for the compact convex
set case) is the most commonly used. A more general version states:
if $C$ is a nonempty convex subset of a Banach space (not necessarily closed
or bounded) and $T \colon C \to C$ is continuous with $T(C)$ relatively
compact, then $T$ has a fixed point. The proof is similar.
\end{remark}

\subsection{Kakutani Fixed Point Theorem}

\begin{definition}[Set-Valued Map / Correspondence]\label{def:set-valued}
A \emph{set-valued map} (or \emph{correspondence}) from a set $X$ to a set $Y$
is a map $\Phi \colon X \rightrightarrows Y$ that assigns to each $x \in X$
a subset $\Phi(x) \subset Y$. A \emph{fixed point} of $\Phi \colon X
\rightrightarrows X$ is a point $x^*$ with $x^* \in \Phi(x^*)$.
\end{definition}

\begin{definition}[Upper Hemicontinuity]\label{def:uhc}
A set-valued map $\Phi \colon X \rightrightarrows Y$ between topological
spaces is \emph{upper hemicontinuous} (u.h.c.) if for every open set
$V \subset Y$, the set $\{x \in X : \Phi(x) \subset V\}$ is open in $X$.
Equivalently (when $Y$ is compact): if $x_n \to x$ and $y_n \in \Phi(x_n)$
with $y_n \to y$, then $y \in \Phi(x)$ (closed graph property when the
image sets are compact).
\end{definition}

\begin{theorem}[Kakutani Fixed Point Theorem]\label{thm:kakutani}
Let $C \subset \R^n$ be a nonempty compact convex set and
$\Phi \colon C \rightrightarrows C$ a set-valued map such that:
\begin{enumerate}[label=(\roman*)]
    \item $\Phi(x)$ is nonempty and convex for every $x \in C$.
    \item $\Phi$ is upper hemicontinuous.
\end{enumerate}
Then $\Phi$ has a fixed point: there exists $x^* \in C$ with
$x^* \in \Phi(x^*)$.
\end{theorem}

\begin{proof}
We prove this by approximation with single-valued maps and application of
Brouwer's theorem.

\textbf{Step 1: Approximate selections.}
For each $k \in \N$, we construct a continuous function
$f_k \colon C \to C$ such that the graph of $f_k$ is ``close'' to the
graph of $\Phi$.

For each $k$, cover $C$ by finitely many open balls $B(x_i, 1/k)$,
$i = 1, \ldots, m_k$, with $x_i \in C$. Let $\{\lambda_i\}$ be a partition
of unity subordinate to this cover. For each $i$, pick any
$y_i \in \Phi(x_i)$. Define
\[
    f_k(x) = \sum_{i=1}^{m_k} \lambda_i(x) \, y_i.
\]
This is continuous, and $f_k(x) \in C$ since $y_i \in C$ and $C$ is convex.

\textbf{Step 2: Apply Brouwer.}
By Brouwer's Fixed Point Theorem, each $f_k$ has a fixed point
$x_k \in C$ with $f_k(x_k) = x_k$.

\textbf{Step 3: Pass to the limit.}
Since $C$ is compact, $(x_k)$ has a convergent subsequence $x_{k_j} \to x^*$.
We claim $x^* \in \Phi(x^*)$.

Fix $\varepsilon > 0$. Since $\Phi$ is u.h.c.\ and $\Phi(x^*)$ is convex
(hence, being a convex compact subset of $\R^n$, it equals the intersection
of all open convex sets containing it), let $V$ be an open convex set with
$\Phi(x^*) \subset V$. By upper hemicontinuity, $\{x : \Phi(x) \subset V\}$
is open and contains $x^*$, so for large $j$, $\Phi(x_{k_j}) \subset V$
and $\|x_{k_j} - x^*\| < \varepsilon$.

For large $k_j > 1/\varepsilon$: $\lambda_i(x_{k_j}) > 0$ only when
$\|x_{k_j} - x_i\| < 1/k_j < \varepsilon$, so $\|x_i - x^*\| < 2\varepsilon$.
For such $i$, when $j$ is large enough, $x_i$ is close to $x^*$, so
$\Phi(x_i) \subset V$ (by upper hemicontinuity), hence $y_i \in V$.
Then $f_{k_j}(x_{k_j}) = x_{k_j}$ is a convex combination of $y_i$'s in $V$.
Since $V$ is convex, $x_{k_j} \in V$.

Since $V$ was an arbitrary open convex neighborhood of $\Phi(x^*)$, and
$x_{k_j} \to x^*$, we get $x^* \in \overline{V}$ for every such $V$.
Since $\Phi(x^*)$ is compact and convex,
$x^* \in \bigcap \{\overline{V} : V \supset \Phi(x^*), V \text{ open convex}\}
= \Phi(x^*)$.
\end{proof}

\begin{rlconnection}[Nash Equilibrium Existence via Kakutani]
Consider a finite $n$-player game. Player $i$ has a finite action set $A_i$
and a payoff function $u_i \colon \prod_j A_j \to \R$. A \emph{mixed strategy}
for player $i$ is a probability distribution $\sigma_i \in \Delta(A_i)$
(the simplex over $A_i$).

The \textbf{best-response correspondence} of player $i$ is
\[
    BR_i(\sigma_{-i}) = \arg\max_{\sigma_i \in \Delta(A_i)}
    \sum_{a \in A} \Bigl(\prod_j \sigma_j(a_j)\Bigr) u_i(a).
\]
Define the joint best-response $\Phi(\sigma) = \prod_i BR_i(\sigma_{-i})$
on $C = \prod_i \Delta(A_i)$, which is a nonempty compact convex set in
$\R^{\sum_i |A_i|}$.

One verifies:
\begin{enumerate}[leftmargin=2em]
    \item Each $BR_i(\sigma_{-i})$ is nonempty and convex (it is the set of
          maximizers of a linear function over a simplex).
    \item $\Phi$ is upper hemicontinuous (by Berge's Maximum Theorem: the
          payoff is continuous and the constraint set is constant).
\end{enumerate}

By Kakutani's theorem, $\Phi$ has a fixed point $\sigma^*$, which is
precisely a \textbf{Nash equilibrium}: every player is best-responding to the
others. This is the original proof by Nash (1950).
\end{rlconnection}

\begin{remark}\label{rem:fixed-point-hierarchy}
The fixed point theorems form a hierarchy:
\[
    \text{Banach} \subset \text{Brouwer} \subset \text{Schauder}
    \quad \text{and} \quad
    \text{Brouwer} \subset \text{Kakutani}.
\]
Banach gives uniqueness and algorithms (iterative computation);
Brouwer/Schauder/Kakutani give only existence. The choice depends on the
structure available:
\begin{center}
\begin{tabular}{lll}
\toprule
\textbf{Theorem} & \textbf{Setting} & \textbf{Gives} \\
\midrule
Banach & Complete metric, contraction & Unique FP + algorithm \\
Brouwer & Compact convex $\subset \R^n$, continuous & Existence \\
Schauder & Bounded closed convex, compact map & Existence \\
Kakutani & Compact convex $\subset \R^n$, u.h.c.\ correspondence & Existence \\
\bottomrule
\end{tabular}
\end{center}
\end{remark}


%% ============================================================================
%% SECTION 9: EXERCISES
%% ============================================================================
\section{Exercises}\label{sec:exercises}

\begin{exercise}[Parallelogram Law Characterization]\label{ex:parallelogram}
Let $(X, \|\cdot\|)$ be a normed space satisfying the parallelogram law:
$\|x + y\|^2 + \|x - y\|^2 = 2(\|x\|^2 + \|y\|^2)$ for all $x, y \in X$.
Show that the formula
\[
    \langle x, y \rangle = \frac{1}{4}(\|x+y\|^2 - \|x-y\|^2)
\]
defines an inner product on $X$ whose induced norm equals $\|\cdot\|$.
\emph{Hint:} Verify additivity $\langle x + z, y \rangle = \langle x, y \rangle
+ \langle z, y \rangle$ first for $z = x$, then for $z = y$, extend to
rational scalars, then use continuity.
\end{exercise}

\begin{exercise}[Non-Inner-Product Spaces]\label{ex:non-ip}
Show that $\ell^1$ and $\ell^\infty$ do not satisfy the parallelogram law,
and therefore are not inner product spaces. Find explicit vectors $x, y$
witnessing the failure.
\end{exercise}

\begin{exercise}[Orthogonal Projection Computation]\label{ex:proj-computation}
Let $H = L^2([0, 1])$ and $M = \spn\{1, t, t^2\}$ (polynomials of degree
$\leq 2$). Compute the orthogonal projection of $f(t) = e^t$ onto $M$ using:
\begin{enumerate}[label=(\alph*)]
    \item The Gram--Schmidt process applied to $\{1, t, t^2\}$.
    \item The normal equations $\langle f - Pf, g \rangle = 0$ for all $g \in M$.
\end{enumerate}
\end{exercise}

\begin{exercise}[Riesz Representation and Gradient]\label{ex:riesz-gradient}
Let $f \colon \R^n \to \R$ be differentiable. The derivative $Df(x)$ at a
point $x$ is a bounded linear functional on $\R^n$. Use the Riesz
Representation Theorem to show that there exists a unique $\nabla f(x) \in \R^n$
such that $Df(x)[h] = \langle \nabla f(x), h \rangle$ for all $h \in \R^n$.
This justifies calling $\nabla f(x)$ ``the gradient.'' How does the gradient
depend on the choice of inner product?
\end{exercise}

\begin{exercise}[Banach Contraction and Successive Approximation]\label{ex:banach-ode}
Consider the integral equation $x(t) = \int_0^t K(t,s) x(s)\, ds + g(t)$ on
$C[0, 1]$ with $\|K\|_\infty \leq M$. Show that the operator
$T(x)(t) = \int_0^t K(t,s) x(s)\, ds + g(t)$ may not be a contraction in
$\|\cdot\|_\infty$, but $T^n$ is a contraction for large $n$, and deduce
that $T$ has a unique fixed point. \emph{Hint:} Show
$\|T^n x - T^n y\|_\infty \leq \frac{(Mt)^n}{n!} \|x - y\|_\infty$
and use the fact that a map with a contractive iterate has a unique fixed point.
\end{exercise}

\begin{exercise}[Bellman Operator is a Contraction]\label{ex:bellman-contraction}
Let $\mathcal{S}$ and $\mathcal{A}$ be finite sets, $R \colon \mathcal{S}
\times \mathcal{A} \to \R$ a reward function, $P(\cdot | s, a)$ a transition
kernel, and $\gamma \in (0, 1)$ a discount factor. Define the Bellman
optimality operator $T$ on $V(\mathcal{S})$ (bounded real-valued functions on
$\mathcal{S}$ with the sup norm) by
\[
    (TV)(s) = \max_{a \in \mathcal{A}}\Bigl[R(s,a) + \gamma \sum_{s'} P(s'|s,a) V(s')\Bigr].
\]
\begin{enumerate}[label=(\alph*)]
    \item Prove that $T$ is a $\gamma$-contraction:
          $\|TV - TW\|_\infty \leq \gamma \|V - W\|_\infty$.
          \emph{Hint:} For each $s$, let $a^*$ achieve the max for $TV(s)$;
          then bound $TV(s) - TW(s) \leq \gamma \sum_{s'} P(s'|s,a^*)|V(s') - W(s')|$.
    \item Deduce from the Banach contraction theorem that value iteration
          converges geometrically and that $V^*$ is the unique solution of
          the Bellman equation.
    \item How many iterations are needed to guarantee
          $\|V_n - V^*\|_\infty \leq \varepsilon$?
\end{enumerate}
\end{exercise}

\begin{exercise}[Closed Graph Theorem Application]\label{ex:closed-graph-app}
Let $H$ be a Hilbert space and $T \colon H \to H$ a linear operator satisfying
$\langle Tx, y \rangle = \langle x, Ty \rangle$ for all $x, y \in H$.
Use the Closed Graph Theorem to show that $T$ is bounded.
\emph{Hint:} Show that $T$ has a closed graph by using the symmetry condition
to identify the limit.
\end{exercise}

\begin{exercise}[Uniform Boundedness Application]\label{ex:ubp-app}
Let $(f_n)$ be a sequence in $H^*$ (the dual of a Hilbert space $H$) such
that $\sum_{n=1}^\infty |f_n(x)|^2 < \infty$ for every $x \in H$. Show that
$\sup_n \|f_n\| < \infty$.
\emph{Hint:} The partial sums $S_N(x) = \sum_{n=1}^N |f_n(x)|^2$ are
pointwise bounded; apply the Uniform Boundedness Principle to appropriate
operators.
\end{exercise}

\begin{exercise}[Kakutani and Two-Player Zero-Sum Games]\label{ex:kakutani-zerosum}
Let $A \in \R^{m \times n}$ be a payoff matrix for a two-player zero-sum game.
Player 1 chooses $p \in \Delta_m$ (the $m$-simplex) and Player 2 chooses
$q \in \Delta_n$, with payoff $p^\top A q$ to Player 1.
\begin{enumerate}[label=(\alph*)]
    \item Define the best-response correspondences $BR_1(q)$ and $BR_2(p)$.
    \item Verify the hypotheses of Kakutani's theorem for the joint
          best-response $\Phi(p,q) = BR_1(q) \times BR_2(p)$ on
          $\Delta_m \times \Delta_n$.
    \item Conclude that a Nash equilibrium (saddle point) exists and relate
          this to the minimax theorem.
\end{enumerate}
\end{exercise}

\begin{exercise}[Schauder Theorem Application]\label{ex:schauder-app}
Let $K \colon [0,1] \times [0,1] \to \R$ be continuous and define
$T \colon C[0,1] \to C[0,1]$ by $(Tf)(t) = \int_0^1 K(t,s) f(s)\, ds$.
\begin{enumerate}[label=(\alph*)]
    \item Show that $T$ maps the closed ball
          $\overline{B}(0, R) \subset C[0,1]$ into itself for
          $R \geq \|K\|_\infty \cdot R$ (choose $R$ appropriately).
    \item Show that $T(\overline{B}(0, R))$ is equicontinuous (use uniform
          continuity of $K$).
    \item Apply the Arzel\`{a}--Ascoli theorem to conclude $T$ is compact.
    \item Apply Schauder's theorem to find a fixed point.
\end{enumerate}
\end{exercise}


%% ============================================================================
%% SECTION 10: SUMMARY
%% ============================================================================
\section{Summary of Key Results}\label{sec:summary}

\begin{center}
\small
\begin{tabular}{p{4.0cm} p{5.5cm} p{4.5cm}}
\toprule
\textbf{Result} & \textbf{Statement (Informal)} & \textbf{RL/ML Connection} \\
\midrule
Cauchy--Schwarz (Thm.~\ref{thm:cauchy-schwarz})
    & $|\langle x,y \rangle| \leq \|x\|\|y\|$
    & Bounds on correlations; kernel evaluations \\
\addlinespace
Parallelogram Law (Thm.~\ref{thm:parallelogram})
    & $\|x+y\|^2 + \|x-y\|^2 = 2(\|x\|^2 + \|y\|^2)$
    & Characterizes inner product spaces \\
\addlinespace
Projection Theorem (Thm.~\ref{thm:projection})
    & $H = M \oplus M^\perp$ for closed $M$
    & LSTD projection; least-squares approximation \\
\addlinespace
Bessel's Inequality (Thm.~\ref{thm:bessel})
    & $\sum |\langle x, e_n \rangle|^2 \leq \|x\|^2$
    & Energy bounds for Fourier coefficients \\
\addlinespace
Parseval's Identity (Thm.~\ref{thm:parseval})
    & $\|x\|^2 = \sum |\langle x, e_n \rangle|^2$ for ONB
    & Isometry between signal and coefficient spaces \\
\addlinespace
Riesz Representation (Thm.~\ref{thm:riesz})
    & Every $f \in H^*$ is $\langle \cdot, y \rangle$
    & RKHS theory; gradient = Riesz representative \\
\addlinespace
Hahn--Banach (Thm.~\ref{thm:hahn-banach})
    & Linear functionals extend with norm preserved
    & Dual characterizations; separation theorems \\
\addlinespace
Uniform Boundedness (Thm.~\ref{thm:ubp})
    & Pointwise bounded $\Rightarrow$ uniformly bounded
    & Convergence of operator sequences \\
\addlinespace
Open Mapping (Thm.~\ref{thm:open-mapping})
    & Surjective bounded linear maps are open
    & Bounded inverse theorem \\
\addlinespace
Closed Graph (Thm.~\ref{thm:closed-graph})
    & Closed graph $\Leftrightarrow$ bounded (Banach spaces)
    & Proving operator continuity \\
\addlinespace
Banach Contraction (Thm.~\ref{thm:banach-contraction})
    & Contraction has unique FP; iterates converge at rate $\gamma^n$
    & Value iteration; policy evaluation; $Q$-learning \\
\addlinespace
Brouwer FPT (Thm.~\ref{thm:brouwer})
    & Continuous self-map of compact convex set has FP
    & Existence results in finite dimensions \\
\addlinespace
Schauder FPT (Thm.~\ref{thm:schauder})
    & Compact self-map of closed bounded convex set has FP
    & Existence of solutions to integral equations \\
\addlinespace
Kakutani FPT (Thm.~\ref{thm:kakutani})
    & u.h.c.\ convex-valued correspondence has FP
    & Nash equilibrium existence \\
\bottomrule
\end{tabular}
\end{center}

\bigskip

\noindent\textbf{Key Takeaways:}
\begin{enumerate}[leftmargin=2em]
    \item \textbf{Inner products provide geometry.} The Cauchy--Schwarz
          inequality, projections, and orthogonal decompositions are all
          consequences of the inner product structure. This geometry is
          what makes Hilbert spaces the natural setting for least-squares
          and projection methods.
    \item \textbf{Hilbert spaces are self-dual.} The Riesz Representation
          Theorem identifies $H$ with $H^*$, vastly simplifying the theory
          compared to general Banach spaces.
    \item \textbf{The four fundamental theorems are structural results} about
          Banach spaces. They guarantee that linear operators and functionals
          behave well: they can be extended (Hahn--Banach), pointwise bounds
          imply uniform bounds (UBP), surjections are open (OMT), and
          closed-graph operators are bounded (CGT).
    \item \textbf{Fixed point theorems bridge analysis and algorithms.}
          The Banach contraction theorem provides both existence, uniqueness,
          and an algorithm. Brouwer/Schauder/Kakutani provide existence in
          settings where contractions are not available, notably for Nash
          equilibria.
    \item \textbf{The hierarchy of fixed point theorems} matches the hierarchy
          of RL/ML problems: contraction for value iteration, Kakutani for
          Nash equilibria, Schauder for existence of solutions to
          functional equations.
\end{enumerate}

\end{document}
